%%=============================================================================
%% Probleemdomein
%%=============================================================================
\chapter{Probleemdomein}%
\label{ch:probleemdomein}

Om een duidelijk inzicht te krijgen in het huidige evaluatieproces van studentenprojecten binnen de opleidingsonderdelen Front-End Web Development en Web Services, werden interviews afgenomen met de drie lectoren die deze projecten evalueren. Het doel van deze interviews was om de huidige werkwijze, de tijdsbesteding, de gebruikte evaluatiecriteria en de ervaren moeilijkheden tijdens het beoordelen van projecten in kaart te brengen.

\section{Tijdsbesteding bij het evalueren van projecten}

Uit de interviews blijkt dat het beoordelen van één studentenproject gemiddeld ongeveer één uur in beslag neemt. Deze tijd kan echter sterk variëren afhankelijk van de kwaliteit en/of complexiteit van het project. Voor grotere of complexere projecten kan de evaluatietijd oplopen tot anderhalf uur, terwijl eenvoudiger projecten sneller beoordeeld kunnen worden. Dit vooral doordat er extra onderzoek gedaan moet worden naar extra technologieën die gebruikt werden en omdat er bij grotere projecten vaak meerdere extra technologieën aanwezig zijn. Projecten die niet correct opstarten of onvoldoende gedocumenteerd zijn vereisen vaak meer tijd omdat de lector eerst moet achterhalen hoe het project werkt.

\vspace{1em}

De evaluatie begint meestal met het binnenhalen van het project, het installeren van dependencies en het opstarten van de applicatie aan de hand van de instructies die de student heeft voorzien in de README. Indien dit proces niet vlot verloopt, kan dit de totale evaluatietijd aanzienlijk verhogen. Daarnaast werd aangegeven dat zowel zeer goede als zeer zwakke projecten vaak meer tijd vragen dan gemiddelde projecten. Bij sterke projecten wordt meer tijd besteed aan het begrijpen van de gebruikte technieken, terwijl bij zwakke projecten extra aandacht besteed wordt aan het formuleren van correcte en duidelijke feedback. Dit wordt gedaan om de studenten die naar tweede zit worden doorverwezen extra duidelijkheid te geven en te helpen om hun project voldoende te kunnen verbeteren voor de volgende evaluatie.

\vspace{1em}

De meerderheid van de lectoren geven aan dat de evaluaties vaak in lange sessies gebeuren, waarbij meerdere projecten na elkaar beoordeeld worden. Hierdoor kan de snelheid van evalueren variëren doorheen de dag. In het begin van een evaluatieperiode wordt vaak trager gewerkt omdat er nog geen duidelijke referentie bestaat voor het niveau van de projecten. Naarmate meer projecten beoordeeld worden, ontstaat een soort referentiekader, waardoor latere evaluaties sneller verlopen. Soms wordt ook teruggekeerd naar eerder beoordeelde projecten om de beoordeling beter te kunnen vergelijken en de consistentie te bewaren. In bijzondere gevallen worden bij twijfel de collega's ingeschakeld, om de evaluatie correct en consistent te houden.

\section{Verdeling tussen controle van criteria en inhoudelijke feedback}

Tijdens de interviews werd ook gevraagd naar welk deel van de evaluatietijd besteed wordt aan het controleren van evaluatiecriteria, en welk deel aan het formuleren van inhoudelijke feedback. De lectoren geven aan dat het controleren van basiscriteria relatief snel kan verlopen, maar dat het opzoeken van de implementatie in de code toch een aanzienlijke hoeveelheid tijd vraagt. Ook niet iedereen gaat op dezelfde manier te werk. Een deel van de lectoren overloopt de lijst op voorhand, een ander deel bekijkt de lijst tijdens het evalueren. Vooral bij grotere projecten kan het zoeken naar specifieke onderdelen in de codebase zeer tijdrovend zijn.

\vspace{1em}

Het grootste deel van de evaluatietijd wordt besteed aan het formuleren van inhoudelijke feedback. Lectoren proberen niet alleen na te gaan of een criterium aanwezig is, maar ook hoe goed het geïmplementeerd is en of de student de onderliggende concepten begrijpt. Het uitschrijven van duidelijke en bruikbare feedback wordt als een van de meest tijdsintensieve onderdelen van de evaluatie beschouwd. In sommige gevallen wordt meer dan de helft van de totale evaluatietijd besteed aan het opstellen van deze feedback. Dit omdat er bij projecten die gelijkaardig zijn aan het voorbeeldproject vaak verteld wordt wat er beter was dan het voorbeeld, en wat er slechter was. Bij slechtere projecten wordt er altijd extra aandacht besteed aan duidelijke feedback, om de student zo goed mogelijk op weg te helpen om hun project tegen de volgende evaluatie goed verbeterd te hebben.

\section{Consistentie en subjectiviteit in beoordelingen}

Een belangrijk probleem dat in alle interviews naar voren kwam, is de moeilijkheid om consistent te blijven bij het beoordelen van een groot aantal projecten. De lectoren geven aan dat hun beoordeling beïnvloed kan worden door eerdere projecten die ze bekeken hebben. Na meerdere zwakke projecten kan een gemiddeld project bijvoorbeeld beter lijken dan het in werkelijkheid is, terwijl na een sterk project een gelijkaardig werk strenger beoordeeld kan worden.

\vspace{1em}

Ook de mentale vermoeidheid tijdens lange evaluatiesessies speelt een rol. Sommige lectoren geven aan dat ze minder projecten per dag proberen te beoordelen om de kwaliteit van hun feedback te behouden. Daarnaast wordt aangegeven dat de eigen emotionele toestand of frustratie over slecht uitgewerkte projecten onbewust invloed kan hebben op de beoordeling.

\section{Huidige hulpmiddelen en automatisering}

Momenteel gebeurt de evaluatie van projecten grotendeels manueel. Er bestaan wel hulpmiddelen om projecten automatisch binnen te halen of te controleren of ze kunnen opstarten, maar er wordt geen gebruik gemaakt van een volwaardig autograding- of feedbacksysteem. Enkele lectoren hebben wel geëxperimenteerd met Large Language Models om bepaalde onderdelen van de evaluatie te automatiseren, zoals het controleren van ontvankelijkheidscriteria of het vergelijken van een project met een voorbeeldoplossing.

\vspace{1em}

Deze experimenten toonden aan dat LLM’s zeker potentieel hebben om bepaalde repetitieve controles sneller uit te voeren, maar dat de betrouwbaarheid van de resultaten nog niet voldoende is om ze zonder controle te gebruiken. Er werd bijvoorbeeld vastgesteld dat een model soms vergelijkbare scores geeft aan zowel goede als zwakke projecten, wat wijst op een gebrek aan diepgang in de evaluatie. Daarnaast kan de verwerkingstijd van grote projecten een probleem vormen, vooral wanneer de volledige codebase in één keer geanalyseerd moet worden. Hier komen dan ook context- en tokenlimitaties naar boven. Hierdoor moet de code opgesplitst worden in kleinere delen, wat de complexiteit van een mogelijke automatisering verhoogt.

\section{Belang van inzicht in het werk van de student}

Een ander aspect dat volgens alle geïnterviewde lectoren moeilijk te automatiseren is, is het beoordelen van in welke mate een student de ingediende code zelf begrijpt en geschreven heeft. Om dit te controleren werd voorheen gebruik gemaakt van een mondelinge verdediging. Dit werd voor enkele jaren afgeschaft, maar met de opkomst van LLM's zal dit komend academiejaar terug worden ingelast. Naast het potentiëel van LLM's om projecten te evalueren, zouden ze ook hierbij kunnen helpen. Het afnemen van persoonlijke mondelinge examens vereist ook veel voorbereidend werk, zoals het opstellen van relevante vragen en het doornemen van het project van elke student. LLM's zouden hier ook tijdsbesparend kunnen werken, door een kleine samenvatting van elk project te kunnen opstellen en interessante onderdelen van bepaalde projecten te kunnen highlighten. Hier kan de lector dan dieper op in tijdens het mondelinge examen.

\section{Juridische en ethische overwegingen}

Naast de praktische en technische uitdagingen bij de implementatie van een LLM-gebaseerde evaluatieflow, spelen er ook belangrijke juridische en ethische overwegingen, in het bijzonder op het vlak van data governance, privacy en intellectueel eigendom. Om hierover duidelijkheid te scheppen, werd advies ingewonnen bij de functionaris gegevensbescherming (DPO) en het diensthoofd studentenaangelegenheden van HOGENT.

\vspace{1em}

Ten eerste zijn er de aspecten rond gegevensbescherming en de Algemene Verordening Gegevensbescherming (AVG/GDPR). Het gebruik van studentenprojecten als data voor een AI-systeem impliceert de verwerking van persoonsgegevens. Omdat het systeem gebruikt kan worden in de context van evaluaties en besluitvorming, bevindt deze toepassing zich mogelijk in de hoogrisicozone volgens de Europese AI Act. Dit vereist naleving van verschillende basisbeginselen, zoals rechtmatigheid, transparantie (gebruikers moeten geïnformeerd zijn over de interactie met het AI-systeem), doelbinding en dataminimalisatie. De AI Act legt bovendien een sterke nadruk op zinvol menselijk toezicht gedurende het hele proces. Voor het onderzoek en een eventuele toepassing betekent dit onder meer dat persoonsgegevens idealiter gepseudonimiseerd worden en dat een veilige infrastructuur voor gegevensverwerking noodzakelijk is.

\vspace{1em}

Ten tweede is er het aspect van intellectueel eigendom. Volgens artikel 55 van het Onderwijs- en Examenreglement (OER) van HOGENT blijft de code en het werk dat studenten creëren in het kader van hun opleiding hun intellectueel eigendom. Het gebruik van deze projecten buiten de initiële onderwijscontext waarvoor ze werden ingediend (in dit geval voor analyse door een AI-model in het kader van dit onderzoek) vereist dan ook de expliciete toestemming van de betrokken studenten. Aangezien dit onderzoek wordt uitgevoerd door een student-onderzoeker en niet namens de instelling, kan er geen beroep gedaan worden op eventuele gebruiksrechten van HOGENT.

\vspace{1em}

Tot slot, wat betreft de rol van de lectoren in een AI-ondersteunde evaluatieflow, bepaalt artikel 30 van het OER dat de uiteindelijke evaluatiebevoegdheid en \newline -verantwoordelijkheid onverkort bij de lector blijven berusten. Hoewel AI-output als hulpmiddel kan dienen ter ondersteuning en voor tijdsbesparing, doet dit geen afbreuk aan de persoonlijke verantwoordelijkheid van de evaluator. Er is dus telkens een noodzaak voor controle door de docent, wat ook in lijn is met de vereisten van de AI Act omtrent menselijke betrokkenheid bij geautomatiseerde besluitvorming.

\section{Conclusie van het probleemdomein}

Op basis van de interviews kan geconcludeerd worden dat het huidige evaluatieproces van studentenprojecten tijdsintensief, repetitief en deels subjectief is. Een groot deel van de tijd gaat naar het controleren van criteria en het formuleren van feedback, terwijl de consistentie tussen beoordelingen niet altijd gegarandeerd kan worden. Hoewel er interesse bestaat in het gebruik van automatisering en Large Language Models, worden deze momenteel nog niet structureel ingezet vanwege beperkingen op vlak van betrouwbaarheid, snelheid en contextverwerking.

\vspace{1em}

Daarnaast brengt een eventuele integratie van AI-systemen in het evaluatieproces cruciale juridische en ethische vereisten met zich mee. Het waarborgen van gegevensbescherming, het respecteren van het intellectueel eigendom van de studenten en het behouden van menselijk toezicht bij de beoordeling blijken absolute randvoorwaarden te zijn.

\vspace{1em}

Desondanks kan er worden geconcludeerd dat er duidelijk veel potentieel zit in het bieden van ondersteuning bij deze tijdrovende evaluaties. Naast tijdswinst is er ook mogelijks verhoogde consistentie bij het inzetten van LLM's en zouden LLM's bovendien extra hulp kunnen bieden bij de voorbereiding van mondelinge examens. Dit vormt de aanleiding voor het verdere onderzoek naar de bruikbaarheid en betrouwbaarheid van LLM-gebaseerde evaluatie binnen deze context, waarbij steeds de geldende privacy- en auteursrechten gerespecteerd moeten worden.
